\documentclass[11pt,a4paper]{article}

% Essential packages
\usepackage[utf8]{inputenc}
\usepackage[T1]{fontenc}
\usepackage[english]{babel}

% Page layout
\usepackage[margin=1in]{geometry}
\usepackage{setspace}
\onehalfspacing

% Math packages
\usepackage{amsmath}
\usepackage{amssymb}
\usepackage{amsthm}

% Graphics and figures
\usepackage{graphicx}
\usepackage{float}
\usepackage{subfig}

% Tables
\usepackage{booktabs}
\usepackage{multirow}

% Citations and bibliography
\usepackage{cite}

% Hyperlinks
\usepackage[colorlinks=true,linkcolor=blue,citecolor=blue,urlcolor=blue]{hyperref}

% Code listings (if needed)
\usepackage{listings}
\usepackage{xcolor}

% Algorithms (if needed)
\usepackage{algorithm}
\usepackage{algpseudocode}

% Custom commands
\newcommand{\email}[1]{\texttt{#1}}
\newcommand{\todo}[1]{\textcolor{red}{\textbf{TODO: #1}}}

% Theorem environments
\newtheorem{theorem}{Theorem}
\newtheorem{lemma}[theorem]{Lemma}
\newtheorem{proposition}[theorem]{Proposition}
\newtheorem{corollary}[theorem]{Corollary}
\theoremstyle{definition}
\newtheorem{definition}{Definition}
\newtheorem{example}{Example}
\theoremstyle{remark}
\newtheorem{remark}{Remark}

% Code listing style
\lstset{
    basicstyle=\ttfamily\small,
    breaklines=true,
    frame=single,
    numbers=left,
    numberstyle=\tiny,
    commentstyle=\color{gray},
    keywordstyle=\color{blue},
    stringstyle=\color{red}
}

% Document metadata
\title{Your Paper Title Here}
\author{
    Author Name\thanks{Corresponding author: \email{author@institution.edu}} \\
    \small Department Name \\
    \small Institution Name \\
    \small City, Country
    \and
    Second Author Name \\
    \small Department Name \\
    \small Institution Name \\
    \small City, Country
}
\date{\today}

\begin{document}

\maketitle

\begin{abstract}
This is the abstract of your paper. It should provide a concise summary of the paper's content, including the problem being addressed, the methodology used, key results, and main conclusions. The abstract should be self-contained and typically ranges from 150 to 250 words.
\end{abstract}

\noindent\textbf{Keywords:} keyword1, keyword2, keyword3, keyword4, keyword5

\section{Introduction}
\label{sec:introduction}

The introduction should provide background on the topic, motivate the research problem, and clearly state the objectives of the paper. It should also outline the structure of the remaining sections.

\subsection{Motivation}

Explain why this research is important and what problems it addresses.

\subsection{Contributions}

List the main contributions of this work:
\begin{itemize}
    \item First contribution
    \item Second contribution
    \item Third contribution
\end{itemize}

\subsection{Paper Organization}

The rest of this paper is organized as follows. Section~\ref{sec:related-work} reviews related work. Section~\ref{sec:methodology} describes our methodology. Section~\ref{sec:results} presents experimental results. Section~\ref{sec:discussion} discusses the findings, and Section~\ref{sec:conclusion} concludes the paper.

\section{Related Work}
\label{sec:related-work}

Review relevant literature and explain how your work differs from or builds upon existing research. Cite relevant papers using \cite{example-citation}.

\section{Methodology}
\label{sec:methodology}

Describe your approach in detail. This section should provide enough information for others to reproduce your work.

\subsection{Problem Formulation}

Define the problem formally. You can use mathematical notation:

\begin{equation}
    f(x) = \int_{a}^{b} g(x) \, dx
    \label{eq:example}
\end{equation}

\subsection{Proposed Approach}

Explain your solution approach. You can include algorithms:

\begin{algorithm}
\caption{Example Algorithm}
\label{alg:example}
\begin{algorithmic}[1]
\State \textbf{Input:} Data $D$, parameter $k$
\State \textbf{Output:} Result $R$
\State Initialize $R \leftarrow \emptyset$
\For{$i = 1$ to $k$}
    \State Process data point $d_i$
    \State Update $R$
\EndFor
\State \Return $R$
\end{algorithmic}
\end{algorithm}

\section{Experimental Results}
\label{sec:results}

Present your experimental setup and results.

\subsection{Experimental Setup}

Describe the datasets, evaluation metrics, and baseline methods used.

\subsection{Results}

Present results with tables and figures.

\begin{table}[htbp]
\centering
\caption{Example Results Table}
\label{tab:results}
\begin{tabular}{@{}lcccc@{}}
\toprule
Method & Metric 1 & Metric 2 & Metric 3 & Metric 4 \\
\midrule
Baseline 1 & 0.85 & 0.78 & 0.92 & 0.88 \\
Baseline 2 & 0.87 & 0.80 & 0.90 & 0.86 \\
\textbf{Our Method} & \textbf{0.92} & \textbf{0.85} & \textbf{0.95} & \textbf{0.91} \\
\bottomrule
\end{tabular}
\end{table}

\begin{figure}[htbp]
\centering
% Uncomment and modify when you have an actual figure
% \includegraphics[width=0.8\textwidth]{figures/example-figure.pdf}
\fbox{\parbox{0.8\textwidth}{\centering [Your figure will be placed here]}}
\caption{Example figure caption. Describe what the figure shows and any important details.}
\label{fig:example}
\end{figure}

As shown in Figure~\ref{fig:example} and Table~\ref{tab:results}, our method achieves superior performance.

\section{Discussion}
\label{sec:discussion}

Interpret the results, discuss their implications, and acknowledge limitations of your work.

\subsection{Analysis}

Provide deeper analysis of the results.

\subsection{Limitations}

Discuss any limitations of your approach or experiments.

\section{Conclusion}
\label{sec:conclusion}

Summarize the main findings and contributions of the paper. Discuss potential future work and broader implications.

\subsection{Future Work}

Outline potential directions for future research.

\section*{Acknowledgments}

Acknowledge funding sources, collaborators, and others who contributed to the work.

% Bibliography
\bibliographystyle{plain}
\bibliography{references}

% If you don't have a .bib file yet, you can use the following format:
% \begin{thebibliography}{99}
% \bibitem{example-citation}
% Author, A., Author, B., and Author, C. (2023). 
% Title of the paper. 
% \textit{Journal Name}, 10(2), 123-145.
% \end{thebibliography}

\end{document}
